\section{Síntesis y ampliación}

\subsection{Una cadena de diseño que atraviesa toda la clase}

Esta clase puede condensarse en ocho pasos: primero, formular la tarea como una predicción del futuro condicionada a la acción; distinguir entre translation y disambiguation; elegir la dirección de la fusión; construir datos alineados espacialmente; aprender la distribución condicional con masked modeling; hacer imputation calibrada a partir de modalidades baratas como H\&E; proponer contrafactuales con un modelo unificado; y, por último, actualizar el modelo de mundo mediante experimentos y datos de fallos. Si se omite cualquiera de estos pasos, el sistema puede degradarse de herramienta científica a demostración vistosa.

\begin{importantbox}{Conclusión central}
El valor de un modelo de mundo multimodal no está en «meter más datos en un Transformer más grande», sino en organizar la información complementaria de distintos sensores como predicciones de intervenciones que puedan verificarse. El modelo debe saber cuándo puede traducir, cuándo necesita desambiguar y cuándo la información simplemente no es identificable, y debe incorporar la incertidumbre a la selección de experimentos.
\end{importantbox}

\subsection{La lista de verificación de ingeniería que más vale llevarse}

\begin{enumerate}
  \item \textbf{Tarea:} ¿el objetivo de predicción corresponde a una decisión real, y no solo a algo fácil de optimizar?
  \item \textbf{Partición:} ¿se particiona por paciente, centro y tiempo para evitar fugas espaciales y por lote?
  \item \textbf{Fusión:} ¿la nueva modalidad aporta translation o disambiguation?
  \item \textbf{Faltantes:} ¿durante el entrenamiento se simula la ausencia de modalidades completas y el presupuesto de medición que habrá en el despliegue?
  \item \textbf{Calibración:} ¿se reporta la cobertura de la incertidumbre por gen, por paciente y por dominio?
  \item \textbf{Contrafactuales:} ¿se verifican la dirección y la magnitud del efecto con datos reales de perturbation?
  \item \textbf{Ciclo cerrado:} ¿los experimentos fallidos se registran y se usan para actualizar el modelo?
\end{enumerate}

\subsection{Preguntas abiertas}

Primero, ¿cómo asignar el cómputo de forma dinámica entre los niveles patient--organ--tissue--cell--molecule, en lugar de aplicar una compresión fija? Segundo, ¿cómo distinguir los estados moleculares identificables en H\&E de los que no lo son en absoluto? Tercero, ¿cómo alinear los counterfactual del modelo con intervenciones causales reales? Cuarto, ¿cómo modelar el sesgo de selección cuando faltan datos de tejido sano y de tratamientos aleatorizados? Quinto, ¿cómo debe evaluarse el valor de un scientist agent mediante experimentos prospectivos y no mediante benchmarks de texto?

Estas preguntas apuntan en conjunto a un estándar más riguroso de AI for Science: el modelo no solo debe generar respuestas razonables, sino también exponer el origen de la evidencia, la distancia entre distribuciones, los intervalos de predicción y compromisos experimentales que puedan fallar. El destino de un modelo de mundo no es un mapa de calor más realista, sino un sistema de investigación que la realidad pueda corregir de manera continua.

\subsection{Lecturas adicionales}

\begin{itemize}[itemsep=0.2em,topsep=0.2em]
  \item Página del curso Stanford CS25 V5: \url{https://web.stanford.edu/class/cs25/past/cs25-v5/}
  \item Video oficial de esta clase: \url{https://www.youtube.com/watch?v=8kXIaUM3h1E}
  \item Palabras clave de métodos: CLIP (representación compartida), Masked Autoencoder (aprendizaje por enmascaramiento), cross-attention (fusión dirigida), adaptive LayerNorm (modulación condicional). Al evaluar, conviene buscar primero datos reservados a nivel de paciente, intervenciones reales, calibración y reproducción entre centros.
\end{itemize}

\end{document}
